\documentclass[12pt]{article}
\usepackage[margin=1in]{geometry}
\usepackage{setspace}
\usepackage{hyperref}
\usepackage{amsmath,amssymb}
\usepackage{lmodern}
\renewcommand{\labelitemi}{\ensuremath{\bullet}}

\hypersetup{
    colorlinks=true,
    linkcolor=blue,
    urlcolor=blue,
    citecolor=blue
}

\begin{document}
\singlespacing

\begin{flushright}
Prof. Humberto Bernal \\
Universidad Colegio Mayor de Cundinamarca \\
Doctor en Economía, Universidad de los Andes, Colombia \\
E-mail: \href{mailto:hbernal@uniandes.edu.co}{hbernal@uniandes.edu.co}, \href{mailto:hbernalc@universidadmayor.edu.co}{hbernalc@universidadmayor.edu.co} \\
ORCID: \href{https://orcid.org/0000-0002-4864-1726}{0000-0002-4864-1726} \\
\today
\end{flushright}

\vspace{1.5em}

\begin{flushleft}
Prof. James J.F. Forest, Editor en Jefe \\
\textit{Terrorism and Political Violence} \\
Routledge, Taylor \& Francis
\end{flushleft}

\vspace{1.5em}

\noindent Estimado Profesor Forest,

\vspace{1.5em}

\noindent Es un gran privilegio presentar el manuscrito titulado \textbf{``Identifying Rational Types in Unknown Environments under an Indirect Dynamic Mechanism''} para su consideración de publicación en \textit{Terrorism and Political Violence}.

Dada la destacada reputación internacional y el liderazgo interdisciplinario de \textit{Terrorism and Political Violence} en la publicación de investigaciones rigurosas sobre las dinámicas estratégicas, conductuales y organizacionales de la violencia política, la insurgencia y el terrorismo, consideramos que la revista constituye el espacio ideal para este estudio.

En escenarios de conflicto armado y contraterrorismo---particularmente con organizaciones insurgentes como las Fuerzas Armadas Revolucionarias de Colombia (FARC), el Ejército de Liberación Nacional (ELN) y grupos paramilitares de las Autodefensas Unidas de Colombia (AUC)---las interacciones estratégicas están marcadas por una profunda asimetría de información, operaciones clandestinas y agregación observacional (\textit{observational pooling}). Los mecanismos clásicos de revelación directa colapsan debido a que los actores violentos no estatales poseen incentivos determinantes para ocultar o falsear sus tipos operativos latentes, capacidades y motivaciones estratégicas.

Este artículo formula un novedoso marco de diseño de mecanismos dinámicos indirectos adaptado a las dinámicas estratégicas y conductuales de organizaciones insurgentes bajo agregación observacional. Fundamentado en microdatos del Centro Nacional de Memoria Histórica (CNMH) de Colombia que abarcan más de cuatro décadas de conflicto, el manuscrito presenta cuatro contribuciones teóricas, computacionales y de política pública alineadas directamente con el alcance de \textit{Terrorism and Political Violence}:

\begin{enumerate}
    \item \textbf{Mecanismo Dinámico Indirecto para el Cribado de Insurgencias:} En lugar de depender de esquemas teóricos de comunicación directa inviables en contextos de violencia armada, formalizamos un mecanismo de cribado dinámico indirecto. El Estado (planificador) despliega instrumentos observables de contraterrorismo y antisecuestro, mientras las células insurgentes (agentes) responden óptimamente en función de su presión financiera privada, cohesión ideológica y vulnerabilidad logística.
    
    \item \textbf{Implementación Estocástica mediante el Proceso Mano de Dios--Guadalupe (MDG):} Con el fin de incorporar fricciones tácticas, fallas de inteligencia y contingencias operativas en entornos de conflicto armado de alto riesgo, el proceso MDG modela endógenamente ruido de ejecución bajo restricciones de máxima entropía. Esto asegura soporte pleno en todas las trayectorias de acción, garantizando la consistencia de la actualización bayesiana fuera de la trayectoria y sustentando la implementabilidad en Equilibrio Bayesiano Perfecto (EBP).
    
    \item \textbf{El Teorema de Identificación de Tipos Racionales y Tipología Insurgente:} La contribución teórica central demuestra que, bajo señales observables recurrentemente separadas (duración del cautiverio y riesgos competitivos de desenlace: rescate militar, pago de rescate, fuga o ejecución), las medidas de trayectoria entre tipos insurgentes alternativos resultan mutuamente singulares. Como resultado, las creencias posteriores bayesianas del planificador estatal convergen casi seguramente al verdadero tipo latente de la organización armada (separando facciones delictivas oportunistas de guerrillas ideológicas disciplinadas).
    
    \item \textbf{Diseño de Políticas Contraterroristas ($\alpha$ vs.\ $\gamma$):} Mediante modelos microeconométricos de duración con riesgos competitivos y aproximadores de redes neuronales de aprendizaje profundo para la función de valor, el trabajo evalúa los trade-offs estructurales entre interdicción financiera ($\alpha$, bloqueo de flujos ilícitos y activos de rescate) y presión operativa táctica ($\gamma$, intervenciones militares y policiales). Se caracteriza la frontera de política que rige la supervivencia de los cautivos, la eficiencia extractiva y la duración del cautiverio, demostrando que una presión militar excesiva puede acelerar perversamente ejecuciones letales si no se calibra adecuadamente con medidas de asfixia financiera.
\end{enumerate}

Para asegurar plena transparencia, reproducibilidad y rigor científico:
\begin{itemize}
    \item Las demostraciones matemáticas completas y los lemas asintóticos han sido formalmente verificados en el Asistente Interactivo de Demostraciones \textbf{Lean 4} (Apéndice Suplementario~B).
    \item Los scripts microeconométricos en Stata y las bases de datos analíticas están documentados en el Apéndice Suplementario~A.
    \item Una herramienta interactiva de simulación está disponible públicamente en \url{https://7kkgn4sfscykkfvv5794qw.streamlit.app/} y descrita en el Apéndice Suplementario~C.
    \item Todos los códigos de replicación, bases de datos, demostraciones en Lean~4 y materiales complementarios se encuentran alojados en el repositorio público de GitHub: \url{https://github.com/Donzhumber/TPV.git}.
\end{itemize}

El manuscrito no ha sido publicado previamente ni se encuentra bajo consideración en ninguna otra revista. Conforme a las directrices de revisión por pares doble ciego de \textit{Terrorism and Political Violence}, se proporcionan versiones anónima y con identificación de autor, así como una portada independiente (\textit{Title Page}).

Agradecemos sinceramente su tiempo, consideración y orientación editorial.

\vspace{2.5em}

\begin{flushleft}
Atentamente,

\vspace{2.5em}

\textbf{Prof. Humberto Bernal, Ph.D.} \\
Universidad Colegio Mayor de Cundinamarca \\
Doctor en Economía, Universidad de los Andes, Colombia
\end{flushleft}

\end{document}
